\documentclass[12pt]{article}

\usepackage{amsmath}
\usepackage{amssymb}
\usepackage{geometry}

\geometry{margin=1in}

\title{CSE400 -- Fundamentals of Probability in Computing\\
Lecture 7: Expectation, CDFs, PDFs and Problem Solving}
\author{Lecture Scribe}
\date{}

\begin{document}

\maketitle

\section{Overview of the Lecture}

This lecture discusses important concepts used in probability theory and statistical modeling:

\begin{enumerate}
\item Cumulative Distribution Function (CDF)
\item Probability Density Function (PDF)
\item Expectation of Random Variables
\item Expectation of a Function of a Random Variable
\item Linear Operations with Expectation
\item Moments and Central Moments
\item Variance, Skewness and Kurtosis
\end{enumerate}

These concepts help describe the behavior, average value, and distribution of random variables.

\section{Cumulative Distribution Function (CDF)}

\subsection{Definition}

The \textbf{Cumulative Distribution Function (CDF)} of a random variable $X$ is defined as

\[
F_X(x) = P(X \leq x)
\]

It represents the probability that the random variable takes a value less than or equal to $x$.

\subsection{Properties of the CDF}

\textbf{1. Non-Decreasing Property}

\[
F_X(x_1) \leq F_X(x_2) \quad \text{if } x_1 < x_2
\]

\textbf{2. Limits at Infinity}

\[
\lim_{x \to -\infty} F_X(x) = 0
\]

\[
\lim_{x \to \infty} F_X(x) = 1
\]

\textbf{3. Right Continuity}

\[
\lim_{h \to 0^+} F_X(x+h) = F_X(x)
\]

\textbf{4. Probability Between Two Values}

\[
P(a < X \leq b) = F_X(b) - F_X(a)
\]

\section{Probability Density Function (PDF)}

\subsection{Definition}

For a continuous random variable, the \textbf{Probability Density Function (PDF)} is defined as

\[
f_X(x) = \frac{d}{dx}F_X(x)
\]

\subsection{Relationship Between PDF and CDF}

\[
F_X(x) = \int_{-\infty}^{x} f_X(t)\,dt
\]

Thus,

\begin{itemize}
\item PDF is the derivative of the CDF
\item CDF is the integral of the PDF
\end{itemize}

\subsection{Properties of PDF}

\textbf{1. Non-negativity}

\[
f_X(x) \ge 0
\]

\textbf{2. Total Probability Equals 1}

\[
\int_{-\infty}^{\infty} f_X(x)\,dx = 1
\]

\textbf{3. Probability Between Two Points}

\[
P(a < X < b) = \int_a^b f_X(x)\,dx
\]

\section{Expectation of Random Variables}

\subsection{Definition}

The \textbf{Expectation} (Expected Value) represents the average value of a random variable over many repeated experiments.

\subsection{Expectation of Discrete Random Variable}

\[
E[X] = \sum x_i p_i
\]

\subsection{Expectation of Continuous Random Variable}

\[
E[X] = \int_{-\infty}^{\infty} x f_X(x)\,dx
\]

\section{Expectation of a Function of a Random Variable}

Sometimes we compute the expectation of a function $g(X)$.

\subsection{Discrete Case}

\[
E[g(X)] = \sum g(x_i) P(X = x_i)
\]

\subsection{Continuous Case}

\[
E[g(X)] = \int_{-\infty}^{\infty} g(x) f_X(x)\,dx
\]

\section{Linearity of Expectation}

For constants $a$ and $b$:

\[
E[aX + b] = aE[X] + b
\]

For two random variables $X$ and $Y$:

\[
E[aX + bY] = aE[X] + bE[Y]
\]

Linearity holds even when the random variables are not independent.

\section{Moments of Random Variables}

The \textbf{$n$th moment about the origin} is defined as

\[
E[X^n]
\]

Moments describe the mean, spread, and shape of a distribution.

\section{Central Moments}

Central moments are calculated about the mean:

\[
E[(X-\mu)^n]
\]

where

\[
\mu = E[X]
\]

\section{Variance}

Variance measures the spread of the distribution.

\[
Var(X) = E[(X-\mu)^2]
\]

Equivalent form:

\[
Var(X) = E[X^2] - (E[X])^2
\]

\section{Skewness}

Skewness measures the asymmetry of a distribution.

\begin{itemize}
\item Positive Skewness: Tail extends to the right.
\item Negative Skewness: Tail extends to the left.
\end{itemize}

\section{Kurtosis}

Kurtosis measures the peakedness and heaviness of the tails of a distribution.

\begin{itemize}
\item High kurtosis: heavy tails and sharp peak
\item Low kurtosis: flatter distribution
\end{itemize}

\section{Summary}

Key concepts covered in this lecture include:

\textbf{CDF}

\[
F_X(x) = P(X \le x)
\]

\textbf{PDF}

\[
f_X(x) = \frac{d}{dx}F_X(x)
\]

\textbf{Expectation}

\[
E[X] = \sum x_i p_i
\]

or

\[
E[X] = \int x f_X(x)dx
\]

Additional topics include expectation of functions, linearity of expectation, moments, variance, skewness, and kurtosis.

\end{document}